%% LaTeX2e class for student theses
%% sections/content.tex
%% 
%% Karlsruhe Institute of Technology
%% Institute for Program Structures and Data Organization
%% Chair for Software Design and Quality (SDQ)
%%
%% Dr.-Ing. Erik Burger
%% burger@kit.edu
%%
%% Version 1.6, 2024-06-07

\chapter{Introduction}
\label{ch:Introduction}

%% -------------------
%% | Example content |
%% -------------------

This is the template for Bachelor's and Master's theses at SDQ.
For more information on the formatting of theses at SDQ, please refer to
\url{https://sdq.kastel.kit.edu/wiki/Ausarbeitungshinweise} or to your advisor.

\section{Spacing and indentation}
To separate parts of text in \LaTeX, please use two line breaks in your source code.
They will then be set with correct indentation.
Do \emph{not} use:
\begin{itemize}
  \itemsep0em
  \item \texttt{\textbackslash\textbackslash}
  \item \texttt{\textbackslash parskip}
  \item \texttt{\textbackslash vskip}
\end{itemize} 
or other commands to manually insert spaces, since they break the layout of this template.

\section{Example: Citation}
\label{sec:Introduction:Citation}
This template is based on \texttt{biblatex} and \texttt{biber}, which is preferred over the
outdated Bib\TeX{} software.
Please adjust your build environment if necessary (see
\url{https://sdq.kastel.kit.edu/wiki/BibTeX-Literaturlisten#biblatex.2Fbiber})

A citation: \cite{becker2008a} 

\section{Example: Figures}
\label{sec:Introduction:Figures}
\begin{figure}
\centering
\includegraphics[width=4cm]{logos/sdqlogo}
\caption{SDQ logo}
\label{fig:sdqlogo}
\end{figure}

A reference: The SDQ logo is displayed in \autoref{fig:sdqlogo}. 
(Use \code{\textbackslash autoref\{\}} for easy referencing.) 

\section{Example: Tables}
The \texttt{booktabs} package offers nicely typeset tables, as in \autoref{tab:atable}.

Please use \texttt{\textbackslash small} as the font size for tables.

\label{sec:Introduction:Tables}
\begin{table}
\centering\small
\begin{tabular}{r l}
\toprule
abc & def\\
ghi & jkl\\
\midrule
123 & 456\\
789 & 0AB\\
\bottomrule
\end{tabular}
\caption{A table}
\label{tab:atable}
\end{table}

\section{Example: Formula}
One of the nice things about the Linux Libertine font is that it comes with
a math mode package.
\begin{displaymath}
f(x)=\Omega(g(x))\ (x\rightarrow\infty)\;\Leftrightarrow\;
\limsup_{x \to \infty} \left|\frac{f(x)}{g(x)}\right|> 0
\end{displaymath}

%% --------------------
%% | /Example content |
%% --------------------

Artificial intelligence is rapidly changing many fields of study.
During software development, many artifacts are created, ranging from the actual project code through documentation to a multitude of formal and informal diagrams. The task of trace link recovery is to find instances of the same element across different artifacts. 
My goal is to improve the rate at which trace links are recovered using prompt engineering. Prompt engineering is the process of refining a prompt for the specific use case. This is usually done in a non systemaic way. In this work I will explore a systematic iterative approach to prompt engineering. 

In previous work by Fuchss et al. %%Todo QUote%%
simple handwritten prompts were used to recover trace links between software documentation and architectural diagrams. 